\chapter{Correlation Functions and Conformal Frames}
\label{ch:volume-iii-correlation-functions-conformal-frames}

\section{One- and Two-Point Functions}

Let \(\mathcal O_i\) be scalar primary operators of dimensions
\(\Delta_i\).  Translation invariance makes a one-point function constant:
\[
  \left\langle \mathcal O_i(x)\right\rangle=c_i.
\]
Scale covariance gives
\[
  c_i=\lambda^{-\Delta_i}c_i
  \qquad
  \text{for all }\lambda>0.
\]
Hence \(c_i=0\) unless \(\Delta_i=0\).  In a unitary theory, the scalar
primary with \(\Delta=0\) is the identity representation.

For two scalar primaries, translation and rotation invariance imply
\[
  \left\langle \mathcal O_i(x)\mathcal O_j(y)\right\rangle
  =
  F_{ij}(|x-y|).
\]
Scale covariance gives
\[
  F_{ij}(r)=\frac{C_{ij}}{r^{\Delta_i+\Delta_j}}.
\]
Inversion covariance then requires \(C_{ij}=0\) unless
\(\Delta_i=\Delta_j\).  Therefore
\begin{equation}
  \left\langle \mathcal O_i(x)\mathcal O_j(y)\right\rangle
  =
  \frac{C_{ij}\,\delta_{\Delta_i,\Delta_j}}
       {|x-y|^{2\Delta_i}}.
  \label{eq:scalar-primary-two-point}
\end{equation}
Within a finite-dimensional space of primaries with the same dimension and the
same spin, one may choose a basis that diagonalizes the Hermitian matrix
\(C_{ij}\).

\section{Spin Structures}

The inversion tensor
\begin{equation}
  I_{\mu\nu}(x)
  =
  \delta_{\mu\nu}-2\frac{x_\mu x_\nu}{x^2}
  \label{eq:inversion-tensor}
\end{equation}
is the orthogonal part of the inversion Jacobian.  It is the elementary
building block for two-point functions with vector or tensor indices.  For
the stress tensor,
\begin{equation}
  \left\langle T_{\mu\nu}(x)T_{\rho\sigma}(0)\right\rangle
  =
  \frac{C_T}{|x|^{2D}}
  I_{\mu\nu,\rho\sigma}(x),
  \label{eq:stress-tensor-two-point}
\end{equation}
where
\[
  I_{\mu\nu,\rho\sigma}(x)
  =
  \frac12\bigl(
    I_{\mu\rho}(x)I_{\nu\sigma}(x)
    +
    I_{\mu\sigma}(x)I_{\nu\rho}(x)
  \bigr)
  -
  \frac1D\delta_{\mu\nu}\delta_{\rho\sigma}.
\]
The coefficient \(C_T\) is part of the intrinsic CFT data because the
normalization of \(T_{\mu\nu}\) is fixed by the translation Ward identity.

\begin{figure}[H]
  \centering
  \resizebox{0.90\textwidth}{!}{%
  \begin{tikzpicture}[x=1cm,y=1cm,every node/.style={font=\scriptsize}]
    \tikzset{
      insertion/.style={circle,fill=violet!80!black,inner sep=1.8pt},
      arrow/.style={-{Latex[length=2mm]},line width=0.8pt},
      cylinder/.style={line width=0.85pt}
    }
    \begin{scope}
      \node[font=\small] at (0,2.05) {flat space};
      \draw[thick] (-2.1,-1.3) rectangle (2.1,1.35);
      \node[insertion] at (-0.75,-0.2) {};
      \node[violet!80!black,below left] at (-0.75,-0.2) {\(\mathcal O_j(0)\)};
      \node[insertion] at (0.95,0.45) {};
      \node[violet!80!black,above right] at (0.95,0.45) {\(\mathcal O_i(x)\)};
      \draw[blue!65!black,line width=0.8pt] (-0.75,-0.2)--(0.95,0.45);
      \node[blue!65!black] at (0.05,0.48) {\(|x|\)};
    \end{scope}

    \draw[arrow] (2.65,0.05)--(4.00,0.05);
    \node[align=center] at (3.33,0.58) {radial\\quantization};

    \begin{scope}[shift={(6.35,0)}]
      \node[font=\small] at (0,2.05) {cylinder};
      \draw[cylinder] (-1.15,1.25) ellipse (1.05 and 0.25);
      \draw[cylinder] (-2.20,1.25)--(-2.20,-1.25);
      \draw[cylinder] (-0.10,1.25)--(-0.10,-1.25);
      \draw[cylinder] (-1.15,-1.25) ellipse (1.05 and 0.25);
      \node[violet!80!black] at (-1.15,-0.70) {\(\ket{\mathcal O_j}\)};
      \node[violet!80!black] at (-1.15,0.78) {\(\bra{\mathcal O_i}\)};
      \draw[green!45!black,line width=0.8pt] (-2.20,0.15) arc[start angle=180,end angle=360,x radius=1.05,y radius=0.25];
      \draw[green!45!black,dashed,line width=0.8pt] (-2.20,0.15) arc[start angle=180,end angle=0,x radius=1.05,y radius=0.25];
      \draw[arrow,orange!85!black] (0.38,-1.10)--(0.38,1.10);
      \node[orange!85!black,right] at (0.38,0.05) {\(\tau\)};
    \end{scope}
  \end{tikzpicture}
  }
  \caption{The two-point coefficient is the radial-quantization inner product
  between the state created at the origin and the conjugate insertion at
  infinity.}
  \label{fig:two-point-cylinder-inner-product}
\end{figure}

\section{Three-Point Functions}

For scalar primaries, conformal covariance fixes the three-point function up
to one coefficient:
\begin{equation}
  \left\langle
    \mathcal O_1(x_1)\mathcal O_2(x_2)\mathcal O_3(x_3)
  \right\rangle
  =
  \frac{C_{123}}
  {|x_{12}|^{\Delta_1+\Delta_2-\Delta_3}
   |x_{23}|^{\Delta_2+\Delta_3-\Delta_1}
   |x_{13}|^{\Delta_1+\Delta_3-\Delta_2}},
  \label{eq:scalar-primary-three-point}
\end{equation}
where
\[
  x_{ij}=x_i-x_j.
\]
The constants \(C_{123}\), after fixing the two-point normalization of
primaries, are the scalar three-point OPE coefficients.  With spin, the same
logic gives a finite-dimensional space of allowed tensor structures; each
structure has its own coefficient.

As an important mixed example, the Ward identity fixes the coefficient of the
stress-tensor three-point function with two identical scalar primaries.  For
\(\Delta_1=\Delta_2=\Delta\),
\begin{equation}
  \left\langle
    \mathcal O(x_1)\mathcal O(x_2)T_{\mu\nu}(x_3)
  \right\rangle
  =
  \frac{C_{\mathcal O\mathcal O T}}
  {|x_{13}|^D|x_{23}|^D|x_{12}|^{2\Delta-D}}
  \left(
    \frac{X_\mu X_\nu}{X^2}
    -
    \frac1D\delta_{\mu\nu}
  \right),
  \label{eq:scalar-scalar-stress-tensor-three-point}
\end{equation}
with
\[
  X^\mu=\frac{x_{13}^\mu}{x_{13}^2}
       -\frac{x_{23}^\mu}{x_{23}^2}.
\]
The constant \(C_{\mathcal O\mathcal O T}\) is determined by
\(\Delta\) and by the two-point normalization of \(\mathcal O\), because
\(T_{\mu\nu}\) generates translations and dilatations.

\section{Four-Point Functions and Cross Ratios}

For four scalar insertions, conformal transformations can remove all position
dependence except two cross ratios:
\begin{equation}
  u=
  \frac{x_{12}^2x_{34}^2}{x_{13}^2x_{24}^2},
  \qquad
  v=
  \frac{x_{14}^2x_{23}^2}{x_{13}^2x_{24}^2}.
  \label{eq:cross-ratios-u-v}
\end{equation}
Equivalently, the four points can be moved into a two-plane and placed at
\[
  z_1=0,\qquad z_2=z,\qquad z_3=1,\qquad z_4=\infty,
\]
with
\[
  u=z\bar z,\qquad v=(1-z)(1-\bar z).
\]

\begin{figure}[H]
  \centering
  \resizebox{0.94\textwidth}{!}{%
  \begin{tikzpicture}[x=1cm,y=1cm,every node/.style={font=\scriptsize}]
    \tikzset{
      insertion/.style={circle,fill=violet!80!black,inner sep=1.8pt},
      arrow/.style={-{Latex[length=2mm]},line width=0.8pt}
    }
    \begin{scope}
      \node[font=\small] at (0,2.0) {four points in \(\R^D\)};
      \draw[thick] (-2.25,-1.35) rectangle (2.25,1.35);
      \foreach \x/\y/\lab in {
        -1.55/-0.70/{x_1},
        -0.55/0.85/{x_2},
         0.85/0.55/{x_3},
         1.45/-0.55/{x_4}
      }{
        \node[insertion] at (\x,\y) {};
        \node[violet!80!black,above right] at (\x,\y) {\(\lab\)};
      }
      \draw[blue!65!black,line width=0.7pt] (-1.55,-0.70)--(-0.55,0.85);
      \draw[blue!65!black,line width=0.7pt] (0.85,0.55)--(1.45,-0.55);
      \draw[orange!85!black,line width=0.7pt] (-1.55,-0.70)--(0.85,0.55);
      \draw[orange!85!black,line width=0.7pt] (-0.55,0.85)--(1.45,-0.55);
    \end{scope}

    \draw[arrow] (2.80,0)--(4.15,0);
    \node[align=center] at (3.47,0.52) {choose a\\conformal frame};

    \begin{scope}[shift={(6.80,0)}]
      \node[font=\small] at (0,2.0) {two-plane frame};
      \draw[thick] (-2.15,-1.15)--(2.15,-1.15);
      \draw[thick] (-1.65,-1.35)--(-1.65,1.35);
      \node[insertion] at (-1.65,-1.15) {};
      \node[violet!80!black,below] at (-1.65,-1.15) {\(0\)};
      \node[insertion] at (-0.55,0.45) {};
      \node[violet!80!black,above] at (-0.55,0.45) {\(z\)};
      \node[insertion] at (0.75,-1.15) {};
      \node[violet!80!black,below] at (0.75,-1.15) {\(1\)};
      \draw[arrow,violet!80!black] (1.25,-1.15)--(1.95,-1.15);
      \node[violet!80!black,below] at (1.95,-1.15) {\(\infty\)};
      \draw[blue!65!black,line width=0.7pt] (-1.65,-1.15)--(-0.55,0.45);
      \draw[orange!85!black,line width=0.7pt] (-1.65,-1.15)--(0.75,-1.15);
      \node[blue!65!black] at (-0.95,-0.25) {\(u=z\bar z\)};
      \node[orange!85!black] at (0.45,-0.55) {\(v=(1-z)(1-\bar z)\)};
    \end{scope}
  \end{tikzpicture}
  }
  \caption{A scalar four-point function has two conformal invariants.  A
  conformal frame places the points at \(0,z,1,\infty\) in a two-plane.}
  \label{fig:four-point-conformal-frame}
\end{figure}

The general scalar four-point function can be written as
\begin{align}
  &\left\langle
    \mathcal O_1(x_1)\mathcal O_2(x_2)
    \mathcal O_3(x_3)\mathcal O_4(x_4)
  \right\rangle
  \notag\\
  &\quad =
  \frac{\mathcal G(u,v)}
  {|x_{12}|^{\Delta_1+\Delta_2}
   |x_{34}|^{\Delta_3+\Delta_4}}
  \left(\frac{|x_{24}|}{|x_{14}|}\right)^{\Delta_1-\Delta_2}
  \left(\frac{|x_{14}|}{|x_{13}|}\right)^{\Delta_3-\Delta_4}.
  \label{eq:scalar-four-point-general}
\end{align}
For identical scalar primaries of dimension \(\Delta\), this simplifies to
\begin{equation}
  \left\langle
    \mathcal O(x_1)\mathcal O(x_2)
    \mathcal O(x_3)\mathcal O(x_4)
  \right\rangle
  =
  \frac{\mathcal G(u,v)}
       {|x_{12}|^{2\Delta}|x_{34}|^{2\Delta}}.
  \label{eq:identical-scalar-four-point}
\end{equation}
The function \(\mathcal G(u,v)\) is dynamical CFT data.  The next stage is to
determine how the operator product expansion decomposes this function into
conformal blocks and how crossing symmetry constrains the allowed spectrum and
OPE coefficients.
